\oldpage[I]{87}
que $s_i$ est la restriction à $D(f_i)$ de $\theta_1(z)$, ce qui prouve que
$s=\theta_1(z)$ et achève la démonstration.

\begin{corollary}[1.3.8]
\label{I.1.3.8-fr}
Soient $M$, $N$ deux $A$-modules~; l'homomorphisme canonique
$u\mapsto\widetilde{u}$ de $\operatorname{Hom}_A(M,N)$ dans
$\operatorname{Hom}_{\widetilde{A}}(\widetilde{M},\widetilde{N})$ est
bijectif. En particulier, les relations $M=0$ et $\widetilde{M}=0$ sont
équivalentes.
\end{corollary}

Considérons en effet l'homomorphisme canonique $v\mapsto\Gamma(v)$ de
$\operatorname{Hom}_{\widetilde{A}}(\widetilde{M},\widetilde{N})$ dans
$\operatorname{Hom}_{\Gamma(\widetilde{A})}
(\Gamma(\widetilde{M}),\Gamma(\widetilde{N}))$~; ce dernier module s'identifie
canoniquement à $\operatorname{Hom}_A(M,N)$ en vertu du th.
(\hyperref[I.1.3.7-fr]{1.3.7}). Il reste à vérifier que
$u\mapsto\widetilde{u}$ et $v\mapsto\Gamma(v)$ sont réciproques l'un de
l'autre~; or, il est évident que $\Gamma(\widetilde{u})=u$ par définition de
$\widetilde{u}$~; et d'autre part, si on pose $u=\Gamma(v)$ pour
$v\in\operatorname{Hom}_{\widetilde{A}}(\widetilde{M},\widetilde{N})$,
l'application
$w:\Gamma(D(f),\widetilde{M})\to\Gamma(D(f),\widetilde{N})$ déduite
canoniquement de $v$ est telle que le diagramme
\[
  \xymatrix{
    M\ar[r]^u\ar[d]_{\rho_{f,1}} & N\ar[d]^{\rho_{f,1}}\\
    M_f\ar[r]_w & N_f
  }
\]
soit commutatif~; on a donc nécessairement $w=u_f$ pour tout $f\in A$
(\hyperref[0.1.2.4-fr]{0, 1.2.4}), ce qui montre que
$\widetilde{\Gamma(v)}=v$.

\begin{corollary}[1.3.9]
\label{I.1.3.9-fr}
\begin{enumerate}
  \item[\rm(i)] Soit $u$ un homomorphisme d'un $A$-module $M$ dans un
  $A$-module $N$~; alors les faisceaux associés à $\operatorname{Ker}u$,
  $\operatorname{Im}u$, $\operatorname{Coker}u$, sont respectivement
  $\operatorname{Ker}\widetilde{u}$, $\operatorname{Im}\widetilde{u}$,
  $\operatorname{Coker}\widetilde{u}$. En particulier pour que
  $\widetilde{u}$ soit injectif (resp. surjectif, bijectif), il faut et il
  suffit que $u$ le soit.
  \item[\rm(ii)] Si $M$ est une limite inductive (resp. somme directe) d'une
  famille de $A$-modules $(M_\lambda)$, $\widetilde{M}$ est limite inductive
  (resp. somme directe) de la famille $(\widetilde{M_\lambda})$, à un
  isomorphisme canonique près.
\end{enumerate}
\end{corollary}

(i) Il suffit d'appliquer le fait que $\widetilde{M}$ est un foncteur exact en
$M$ (\hyperref[I.1.3.5-fr]{1.3.5}) aux deux suites exactes de $A$-modules
\[
  0\to\operatorname{Ker}u\to M\to\operatorname{Im}u\to0
\]
\[
  0\to\operatorname{Im}u\to N\to\operatorname{Coker}u\to0
\]
La seconde assertion résulte alors du th.
(\hyperref[I.1.3.7-fr]{1.3.7}).

(ii) Soit $(M_\lambda,g_{\mu\lambda})$ un système inductif de $A$-modules, de
limite inductive $M$, et soit $g_\lambda$ l'homomorphisme canonique
$M_\lambda\to M$. Comme on a
$\widetilde{g_{\nu\mu}}\circ\widetilde{g_{\mu\lambda}}
=\widetilde{g_{\nu\lambda}}$ et
$\widetilde{g_\lambda}=\widetilde{g_\mu}\circ
\widetilde{g_{\mu\lambda}}$ pour $\lambda\leq\mu\leq\nu$,
$(\widetilde{M_\lambda},\widetilde{g_{\mu\lambda}})$ est un système
inductif de faisceaux sur $X$, et si on désigne par $h_\lambda$
l'homomorphisme canonique
$\widetilde{M_\lambda}\to\varinjlim\widetilde{M_\lambda}$, il y a un
homomorphisme unique
$v:\varinjlim\widetilde{M_\lambda}\to\widetilde{M}$ tel que
$v\circ h_\lambda=\widetilde{g_\lambda}$. Pour voir que $v$ est bijectif, il
suffit de vérifier que, pour tout $x\in X$, $v_x$ est une bijection de
$(\varinjlim\widetilde{M_\lambda})_x$ sur $\widetilde{M}_x$~; mais
$\widetilde{M}_x=M_x$ et
\[
  (\varinjlim\widetilde{M_\lambda})_x
  =\varinjlim(\widetilde{M_\lambda})_x
  =\varinjlim(M_\lambda)_x=M_x
  \quad(\hyperref[0.1.3.3-fr]{0, 1.3.3}).
\]
D'autre part, il résulte des définitions que $(\widetilde{g_\lambda})_x$ et
$(h_\lambda)_x$ sont tous deux égaux à l'application canonique de
$(M_\lambda)_x$ dans $M_x$~; comme
$(\widetilde{g_\lambda})_x=v_x\circ(h_\lambda)_x$, $v_x$ est l'identité.
